%=========== ΛΥΣΗ - 41ο ΘΕΜΑ Δ ===========
\begin{thema}{Δ}
\begin{erwthma}
\item Θεωρούμε τη συνάρτηση
\[
h(x)=e^x-x,\quad x\in\mathbb{R}.
\]
Για κάθε $x\in\mathbb{R}$ είναι
\[
h'(x)=e^x-1.
\]
Άρα η $h$ είναι γνησίως φθίνουσα στο $(-\infty,0]$ και γνησίως αύξουσα στο $[0,+\infty)$. Επομένως παρουσιάζει ολικό ελάχιστο στο $x_0=0$, με τιμή
\[
h(0)=1.
\]
Συνεπώς
\[
e^x-x\geq1>0
\]
για κάθε $x\in\mathbb{R}$, οπότε η $f$ ορίζεται στο $\mathbb{R}$.

Η $f$ είναι παραγωγίσιμη στο $\mathbb{R}$ και
\[
f'(x)=\frac{e^x-1}{e^x-x}.
\]
Επειδή $e^x-x>0$, το πρόσημο της $f'$ είναι το πρόσημο της παράστασης $e^x-1$. Έχουμε:
\begin{itemize}
\item $f'(x)<0$ για $x<0$,
\item $f'(0)=0$,
\item $f'(x)>0$ για $x>0$.
\end{itemize}
Στον παρακάτω πίνακα βλέπουμε το πρόσημο της $f'$ και τη μονοτονία της $f$:
\begin{center}
\begin{tikzpicture}
\tkzTabInit[color,lgt=2,espcl=1.8,colorC = maincolor!50!white,colorL = maincolor!20!white,colorV = maincolor!50!white]%
{$x$ / .8,$f'(x)$ /0.8,$f(x)$/1}%
{$-\infty$,$0$,$+\infty$}%
\tkzTabLine{, -, z, +, }
\tkzTabVar{+/{+\infty},-/{0},+/{+\infty}}
\end{tikzpicture}
\end{center}

Άρα η $f$ είναι \textbf{γνησίως φθίνουσα} στο $(-\infty,0]$ και \textbf{γνησίως αύξουσα} στο $[0,+\infty)$. Επομένως παρουσιάζει ολικό ελάχιστο στο $x_0=0$, με τιμή
\[
{f(0)=0}.
\]

Από την ανισότητα $e^x-x\geq1$ που αποδείξαμε, προκύπτει
\[
{e^x\geq x+1}
\]
για κάθε $x\in\mathbb{R}$. Η ισότητα ισχύει μόνο για $x=0$.

\item Για κάθε $x\in\mathbb{R}$ είναι
\begin{align*}
f''(x)
&=\left(\frac{e^x-1}{e^x-x}\right)'\\
&=\frac{e^x(e^x-x)-(e^x-1)^2}
{(e^x-x)^2}\\
&=\frac{e^x(2-x)-1}{(e^x-x)^2}.
\end{align*}
Θεωρούμε τη συνάρτηση
\[
g(x)=e^x(2-x)-1.
\]
Έχουμε
\[
g'(x)=e^x(1-x).
\]
Άρα η $g$ είναι γνησίως αύξουσα στο $(-\infty,1]$ και γνησίως φθίνουσα στο $[1,+\infty)$. Επιπλέον,
\[
\lim_{x\to-\infty}g(x)=-1<0,
\quad
g(0)=1>0,
\]
οπότε υπάρχει μοναδική ρίζα
\[
\rho_1\in(-\infty,0).
\]
Επίσης
\[
g(1)=e-1>0
\quad\text{και}\quad
g(2)=-1<0,
\]
οπότε υπάρχει μοναδική ρίζα
\[
\rho_2\in(1,2).
\]

Επειδή $(e^x-x)^2>0$, το πρόσημο της $f''$ είναι το πρόσημο της $g$. Επομένως:
\begin{itemize}
\item $f''(x)<0$ στα διαστήματα $(-\infty,\rho_1)$ και $(\rho_2,+\infty)$,
\item $f''(x)>0$ στο διάστημα $(\rho_1,\rho_2)$.
\end{itemize}
Άρα η $f$ είναι κοίλη στα $(-\infty,\rho_1]$, $[\rho_2,+\infty)$ και κυρτή στο $[\rho_1,\rho_2]$. Η $C_f$ έχει ακριβώς δύο σημεία καμπής
\[
{M_1\left(\rho_1,f(\rho_1)\right)}
\quad\text{και}\quad
{M_2\left(\rho_2,f(\rho_2)\right)},
\]
με $\rho_1<0$ και $\rho_2>1$.

Επειδή
\[
[0,1]\subset(\rho_1,\rho_2),
\]
η $f$ είναι κυρτή στο $[0,1]$.

\item Για την ασύμπτωτη στο $+\infty$ έχουμε
\begin{align*}
\lim_{x\to+\infty}\left[f(x)-x\right]
&=\lim_{x\to+\infty}
\left[\ln(e^x-x)-x\right]\\
&=\lim_{x\to+\infty}
\ln\left(1-xe^{-x}\right)\\
&=0.
\end{align*}
Επομένως η ευθεία
\[
{y=x}
\]
είναι πλάγια ασύμπτωτη της $C_f$ στο $+\infty$.

Επιπλέον,
\[
\lim_{x\to-\infty}f(x)
=\lim_{x\to-\infty}\ln(e^x-x)
=+\infty
\]
και
\[
\lim_{x\to+\infty}f(x)=+\infty.
\]
Επειδή $f(0)=0<2026$, από το Θεώρημα Ενδιάμεσων Τιμών υπάρχει μία τουλάχιστον αρνητική και μία τουλάχιστον θετική ρίζα της εξίσωσης
\[
f(x)=2026.
\]
Η αρνητική ρίζα είναι μοναδική, διότι η $f$ είναι γνησίως φθίνουσα στο $(-\infty,0]$, ενώ η θετική είναι μοναδική, διότι η $f$ είναι γνησίως αύξουσα στο $[0,+\infty)$. Άρα η εξίσωση έχει ακριβώς δύο πραγματικές ρίζες, μία αρνητική και μία θετική.

\item Έχουμε
\[
f(0)=0
\quad\text{και}\quad
f'(0)=0.
\]
Εφαρμόζοντας δύο φορές τον κανόνα \eng{De l'Hospital}, παίρνουμε
\begin{align*}
\lim_{x\to0}\frac{f(x)}{x^2}
&\xlongequal[\text{\eng{DLH}}]{\frac{0}{0}}
\lim_{x\to0}\frac{f'(x)}{2x}\\
&\xlongequal[\text{\eng{DLH}}]{\frac{0}{0}}
\lim_{x\to0}\frac{f''(x)}{2}\\
&=\frac{f''(0)}{2}\\
&={\frac{1}{2}}.
\end{align*}

Η $f$ είναι αυστηρά κυρτή στο $[0,1]$. Επομένως η γραφική της παράσταση βρίσκεται κάτω από τη χορδή που ενώνει τα σημεία
\[
A(0,f(0))
\quad\text{και}\quad
B(1,f(1)).
\]
Επειδή
\[
f(0)=0
\quad\text{και}\quad
f(1)=\ln(e-1),
\]
η εξίσωση της χορδής είναι
\[
y=x\ln(e-1).
\]
Άρα, για κάθε $x\in(0,1)$,
\[
f(x)<x\ln(e-1).
\]
Ολοκληρώνοντας στο $[0,1]$, προκύπτει
\begin{align*}
\int_0^1f(x)\d x
&<\ln(e-1)\int_0^1x\d x\\
&={\frac{\ln(e-1)}{2}}.
\end{align*}
\end{erwthma}
\end{thema}
%=========== ΤΕΛΟΣ ΛΥΣΗΣ - 41ο ΘΕΜΑ Δ ===========
